% Work-MAT | 2. rocnik | pripravovana hodina 21
% Návaznost: uživatelské připravy.zip (hodiny 01–20), autorská učebnice na GitHubu.
% Realisticky.cz: didaktická inspirace v posloupnosti a typech úloh; zadání jsou vlastní.
\documentclass[11pt,a4paper]{article}
\input{../../sablony/pripravy/priprava.sty}
\nastavPripravu{2. ročník}{Exponenciální rovnice}{21}{Lineární substituce}
\begin{document}
\blok{Motivace a návaznost}
Po vytýkání poznáme jiný zápis téže myšlenky: opakující se kladnou mocninu dočasně nahradíme písmenem $t$. Z algebraicky jednoduché rovnice potom vrátíme původní mocninu.
\poznamka{Orientační nabídka pro 45--55 minut: krátké zopakování a motivace, nový výklad, tři společné vzorové úlohy, samostatná práce. Vyber podle tempa třídy; zbývající zadání slouží jako rezerva.}
\blok{Výklad}
Při substituci $t=a^x$ vždy píšeme $t>0$ pro $a>0$, $a\ne1$. Při návratu musí vyhovovat původní exponenciální rovnici.
\blok{Společné příklady}
\begin{prikladReseniBox}[title={Příklad 1}]
Vyřeš, resp. urči: $3\cdot 2^x+2^x=32$.
\tcblower
\[\begin{aligned}t=2^x>0 \\ 3t+t=32\;\Longrightarrow\;t=8 \\ 2^x=8=2^3\;\Longrightarrow\;x=3\end{aligned}\]
\end{prikladReseniBox}
\begin{prikladReseniBox}[title={Příklad 2}]
Vyřeš, resp. urči: $5\cdot3^x-2\cdot3^x=81$.
\tcblower
\[\begin{aligned}t=3^x>0 \\ 3t=81\;\Longrightarrow\;t=27 \\ 3^x=3^3\;\Longrightarrow\;x=3\end{aligned}\]
\end{prikladReseniBox}
\begin{prikladReseniBox}[title={Příklad 3}]
Vyřeš, resp. urči: $4\cdot 5^x+5^x=1$.
\tcblower
\[\begin{aligned}5t=1\;\Longrightarrow\;t=\frac{1}{5}>0 \\ 5^x=5^{-1}\;\Longrightarrow\;x=-1\end{aligned}\]
\end{prikladReseniBox}
\poznamka{Nepřeskoč zápis $t>0$. Po závěru napiš množinu kořenů a zkoušku alespoň u prvního příkladu. Zeptej se, zda by zde nešlo také vytýkat.}
\blok{Samostatné procvičování}
\begin{zadaniReseniBox}[title={Samostatná práce 1}]
\begin{enumerate}[label=\alph*)]
\item $2^x+3\cdot2^x=16$.
\item $3\cdot3^x+6\cdot3^x=81$.
\item $7\cdot2^x-3\cdot2^x=8$.
\item $4\cdot5^x+5^x=25$.
\item $3^x+3^x=\frac{2}{3}$.
\item $2\cdot4^x-4^x=16$.
\end{enumerate}
\tcblower
\begin{enumerate}[label=\alph*)]
\item $x=2$.
\item $x=2$.
\item $x=1$.
\item $x=1$.
\item $x=-1$.
\item $x=2$.
\end{enumerate}
\end{zadaniReseniBox}
\blok{Dobrovolné rozšíření}
\begin{zadaniReseniBox}[title={Rozšíření 1}]
\begin{enumerate}[label=\alph*)]
\item $2^x+2^{x+1}=12$.
\end{enumerate}
\tcblower
\begin{enumerate}[label=\alph*)]
\item $t=2^x>0; 3t=12; x=2.$.
\end{enumerate}
\end{zadaniReseniBox}
\blok{Úkol na doma}
Při početních úlohách napiš postup, zvolenou metodu a zkoušku v původním zadání. U soustav kontroluj obě rovnice; u nerovnic výsledek zakresli na osu.
\begin{zadaniReseniBox}[title={Domácí úkol}]
\begin{enumerate}[label=\alph*)]
\item $2\cdot2^x+2^x=24$.
\item $5\cdot3^x-2\cdot3^x=27$.
\item $3\cdot 4^x+4^x=4$.
\item $7\cdot2^x-5\cdot2^x=\frac{1}{2}$.
\item $3^x+2\cdot3^x=1$.
\item $5^x+4\cdot5^x=125$.
\item $2\cdot7^x-7^x=49$.
\item Vysvětli, proč se nesmí dosadit $t=0$..
\end{enumerate}
\tcblower
\begin{enumerate}[label=\alph*)]
\item $x=3$.
\item $x=2$.
\item $x=0$.
\item $x=-2$.
\item $x=-1$.
\item $x=2$.
\item $x=2$.
\item Protože $a^x>0$..
\end{enumerate}
\end{zadaniReseniBox}
\poznamka{Domácí úkol obsahuje 8 položek; předpokládaná domácí práce přibližně 15--30 minut. Po větších blocích jsou vědomě promíchané dříve probrané typy.}
\end{document}
